\section{Third completion moments with dependent triples included}
Fix an odd prime $p$, length $N=4p$, and a totally isotropic retained
space $G$ of codimension $h\ge2$ in a self-dual code. In the split
$2h$-dimensional quotient $G^\perp/G$, give each isotropic projective
point $\ell$ weight $w_\ell$ equal to the number of normalized sign
classes with nonzero quotient on that point. Let $b$ count the sign
classes in $G$. For a uniformly chosen completion code write
$Y=|S(C)|-b$. Put $D_j=p^{h-j}+1$.

For $t=1,2,3$ let $R_t$ be the sum of
$w_{\ell_1}w_{\ell_2}w_{\ell_3}$ over ordered triples whose span is
totally isotropic of dimension exactly $t$. Repetitions are allowed.
In particular $R_1=\sum_\ell w_\ell^3$. For $h=2$ necessarily
$R_3=0$; the definition of $D_3$ is then unnecessary.

\begin{theorem}[Exact third-moment certificate]
For $h\ge3$,
\[
 \mathbb E Y^3=\frac{R_1}{D_1}+
 \frac{R_2}{D_1D_2}+\frac{R_3}{D_1D_2D_3}.
\]
For $h=2$ omit the final term. If $\mathbb EY^2>0$, some completion
satisfies $|S(C)|\ge b+\lceil\mathbb EY^3/\mathbb EY^2\rceil$.
\end{theorem}
\begin{proof}
The number of generators in a split $2j$-space is
$L_j=2\prod_{i=1}^{j-1}(p^i+1)$ for $j\ge1$, and $L_0=1$.
A specified totally isotropic $t$-space lies in proportion
$L_{h-t}/L_h=\prod_{j=1}^tD_j^{-1}$ of the generators.
This includes $t=h$, where the final factor is $D_h=2$.
A triple not spanning an isotropic space cannot occur in a generator.
Expanding the cube of $Y$ and grouping by span rank proves the formula.
Finally $Y^3\le(\max Y)Y^2$ pointwise gives the certificate.
\end{proof}
Distinct points need not be linearly independent: three different
points on one isotropic projective line contribute to $R_2$, not
$R_3$. Omitting these triples would give an incorrect moment.

\subsection{An exact computation without cubic enumeration}
Let $\mathcal L$ be the totally isotropic projective lines. Then
\[
 R_2=\sum_{L\in\mathcal L}
 \left[\left(\sum_{\ell\in L}w_\ell\right)^3
              -\sum_{\ell\in L}w_\ell^3\right].
\]
Each rank-two triple spans a unique such line, so it is counted once.
Alternatively form the simple orthogonality graph on the weighted
points. Let $T_\triangle$ be the weighted sum over unordered triangles,
and $T_{\rm line}$ the weighted sum over unordered collinear triples.
For unordered orthogonal pairs $\{\ell,m\}$,
\[
 R_2=3\sum_{\{\ell,m\}}w_\ell w_m(w_\ell+w_m)+6T_{\rm line},
 \qquad R_3=6(T_\triangle-T_{\rm line}).
\]
Bitset intersections count graph triangles efficiently when all point
weights are one. For each pair, the other $p-1$ points on its line
are obtained by normalizing $x+ty$, $t\in\mathbb F_p^*$; dictionary
lookup counts collinear triples. Requiring the third index to exceed
both pair indices counts each once. The method uses quadratic pair
enumeration and avoids testing all ordered point triples.

\subsection{Known-order exact benchmarks}
For the first eight stored rows of the order-twenty terminal frame,
$b=9$, $h=2$, and
\[
 R_1=258,\quad R_2=7632,\quad R_3=0,\quad
 \mathbb EY^2=65,\quad\mathbb EY^3=679.
\]
Thus $|S(C)|\ge9+\lceil679/65\rceil=20$.
The second-moment bound gave only nineteen. Independent enumeration
of all twelve generators gives additional counts $3$ six times and
$11$ six times, directly checking the cubic moment.

For the first five stored rows, $b=5$, $h=5$, and the 1092 weighted
points all have weight one. The graph has 6497280 unordered triangles,
of which 1200 are collinear. Exact arithmetic gives
\[
 R_1=1092,\quad R_2=1118736,\quad R_3=38976480,
 \qquad\mathbb EY^3=\frac{995014}{28483}.
\]
As $\mathbb EY^2=14114/2191$, the ratio is
$497507/91741$, giving only $5+\lceil497507/91741\rceil=11$
total classes. This improves the second-moment lower bound of nine,
but does not reach the sufficient threshold seventeen.

\subsection{Scope of the moment hierarchy}
For any integer $r\ge2$ with positive denominator,
\[
 \max Y\ge\frac{\mathbb EY^r}{\mathbb EY^{r-1}}.
\]
These ratios are nondecreasing by Cauchy--Schwarz and tend to
$\max Y$ on the finite completion space. Since $Y$ is integer valued,
their ceilings eventually equal its exact maximum. This is a
classical moment principle, not a uniform bound on the number of
moments required or their computational cost. Supplying sufficiently
large finite-geometry counts remains the missing existence input in
unknown orders. No all-order positivity assertion follows.
